\section{Validation}

We validate the priority corridor rankings through three tests: a coverage simulation that re-runs the 30-mile analysis with each proposed corridor added to the graph, a consistency check against the FHWA Future Interstate Study rankings, and an economic opportunity correlation check.

\subsection{Coverage Simulation}

For each of the 12 priority corridors, we add the proposed alignment to the HighwayGraph as a set of edges with estimated interchange nodes at 25-mile intervals, then re-run the county centroid coverage analysis. Table~\ref{tab:validation} shows the coverage improvement for each corridor.

\begin{table}[h!]
\centering
\caption{Coverage Improvement by Proposed Corridor (Simulation Results)}
\label{tab:validation}
\begin{tabular}{lrrr}
\toprule
\textbf{Corridor} & \textbf{Gap counties served} & \textbf{Gap pop served} & \textbf{US coverage after} \\
\midrule
I-3 (Savannah–Detroit) & 87 & 2,180,000 & 80.3\% \\
Northern Tier (US-2) & 131 & 3,200,000 & 81.5\% \\
I-14 Gulf Coast & 44 & 1,100,000 & 80.0\% \\
I-69 completion & 28 & 750,000 & 79.8\% \\
Northern Maine Interstate & 11 & 830,000 & 79.9\% \\
I-73 (Charlotte–Akron) & 22 & 510,000 & 79.7\% \\
Appalachian Connector & 19 & 620,000 & 79.8\% \\
I-57 N. extension & 7 & 170,000 & 79.6\% \\
Gulf Coast Interior & 27 & 720,000 & 79.8\% \\
I-87 (ND) & 9 & 230,000 & 79.7\% \\
Oregon Interior & 12 & 350,000 & 79.7\% \\
Southwest Connector & 15 & 410,000 & 79.7\% \\
\midrule
\textbf{All 12 combined} & \textbf{412} & \textbf{11,070,000} & \textbf{83.0\%} \\
\bottomrule
\end{tabular}
\end{table}

The combined effect of all 12 corridors raises coverage from 79.6\% to 83.0\% of the US population. This is lower than the first-pass efficiency estimates (which assumed all gap counties within a corridor's 30-mile zone would be served) because some gap counties are at the edge of multiple corridors' service areas and are counted only once in the simulation.

\subsection{Consistency with FHWA Future Interstate Study}

The FHWA Future Interstate Study (2000) \citep{FHWA_FIS2000} identified I-69, I-73, I-3, and the Northern Tier as high-priority proposed corridors based on projected 2020 traffic volumes. Our coverage-efficiency ranking places I-3 (1st) and Northern Tier (2nd) at the top --- consistent with the FHWA study's volume-based priorities. I-69 appears in both studies at comparable priority (4th in our ranking). I-73 appears lower in our ranking than in the FHWA study because its traffic volume justification is stronger than its coverage justification: the Charlotte–Akron corridor serves a moderately dense corridor but not the most geographically isolated populations.

The main divergence: our analysis elevates I-14 (3rd) higher than the FHWA study, which did not prioritize the Gulf South interior. This reflects our coverage-first methodology versus the FHWA's volume-based approach. I-14's relatively low expected traffic volume (it runs through rural Louisiana and Mississippi) makes it less attractive under the volume criterion but highly attractive under the coverage criterion.

\subsection{Economic Opportunity Correlation}

We test whether the priority corridors disproportionately serve counties with high C3 (Economic Opportunity Access) scores --- i.e., whether the corridors with highest coverage efficiency also address the deepest economic gaps.

Mean C3 score of gap counties served by each corridor's top 5 ranked corridors:
\begin{itemize}
\item I-3 served counties: mean C3 = 6.9 (Appalachian coalfield counties --- among highest C3 in dataset)
\item Northern Tier served counties: mean C3 = 5.8 (agricultural northern plains)
\item I-14 served counties: mean C3 = 7.1 (rural Gulf South --- highest C3 zone)
\item I-69 served counties: mean C3 = 6.2 (rural East Texas)
\end{itemize}

All top-4 priority corridors serve counties with C3 scores above the national gap-county mean (5.8), confirming that the coverage-efficiency ranking also selects for economic opportunity impact. The corridors that close the most coverage gaps also serve the populations with the greatest economic opportunity deficit. This alignment --- coverage gap and economic gap pointing to the same investment targets --- strengthens the investment case considerably.

\subsection{Sensitivity to Cost Assumptions}

Our efficiency rankings are sensitive to cost assumptions. We test three cost scenarios:
\begin{itemize}
\item \textbf{Base}: US highway upgrade at \$30M/mi, greenfield at \$75M/mi
\item \textbf{High cost} (+50\%): US highway upgrade at \$45M/mi, greenfield at \$112M/mi
\item \textbf{Low cost} (--30\%): US highway upgrade at \$21M/mi, greenfield at \$53M/mi
\end{itemize}

The rank ordering of the top 5 corridors is stable across all three scenarios. I-3, Northern Tier, and I-14 remain top-3 in all scenarios. The ranking of lower-priority corridors is more sensitive to cost assumptions; we flag these as lower-confidence rankings.
